% !TEX TS-program = xelatex
% !TEX encoding = UTF-8 Unicode
% !BIB program = biber

\documentclass[dvisvgm,aspectratio=169]{beamer}
\usepackage[style=authoryear-icomp, mincitenames=1, maxcitenames=2, url=false, backend=biber]{biblatex}
\usepackage{amsmath, amssymb, mathtools}
\usepackage{textgreek}
\usepackage{caption}
\usepackage{hyperref}
\usepackage[absolute,overlay]{textpos}
\usepackage{appendixnumberbeamer}
\usepackage{pgfpages}
\usepackage{xpatch}
\usepackage{tabularx}
\usepackage{tikz}
\usepackage{animate}

% \usepackage{gfsneohellenicot}

% \usepackage{cmbright}
% \usepackage[OT1]{fontenc}

\usetheme{metropolis}
\usepackage{fontspec}
% \usepackage[sfdefault]{atkinson}
% \setsansfont[BoldFont={Atkinson Hyperlegible}]{Atkinson Hyperlegible}
% \setsansfont[Boldfont={Utopia}]{Utopia}
% \usepackage{noto}
\usepackage{lmodern}

%--------------------------------
% DRAFTING AND SLIDE CONSTRUCTION
%--------------------------------
% \pgfpagesuselayout{4 on 1}
\setbeameroption{hide notes}
% \setbeameroption{show only notes}
% \setbeameroption{show notes on second screen=right}

%-------------------------
% PATCHING & CUSTOMISATION
%-------------------------

\definecolor{UOABlue}{HTML}{00457D}
\definecolor{EngPurple}{HTML}{4F2D7F}
\definecolor{Mustard}{HTML}{906E01}
\setbeamercolor{frametitle}{bg=UOABlue}
\setbeamercolor{palette primary}{bg=UOABlue}
\setbeamercolor{progress bar}{fg=orange}
\metroset{titleformat=smallcaps}
\metroset{sectionpage=simple}
\metroset{progressbar=frametitle}

% \makeatletter
% \setlength{\metropolis@progressinheadfoot@textwidth}{2pt}
% \makeatother

\renewcommand*{\footnotesize}{\tiny}
\renewcommand*{\bibfont}{\scriptsize}
\renewcommand*{\thefootnote}{\fnsymbol{footnote}}

\newcommand{\norm}[2][2]{\left\lVert #2 \right\rVert_{#1}}
\newcommand{\twraps}[2][2cm]{\text{\parbox{#1}{\centering #2}}}
\newcommand{\tPhi}{\text{\textbf{\textPhi}}}

\DeclareMathOperator*{\argmin}{arg\,min}
\DeclareMathOperator*{\argmax}{arg\,max}

\xapptobibmacro{cite}{\setunit{\nametitledelim}\printfield{title}}{}{}

\DeclareSymbolFont{pazosymbols}{OMS}{zplm}{m}{n}
\makeatletter
\ExplSyntaxOn %define \xmapsto
\DeclareMathSymbol \c_pazo_minus \mathbin { pazosymbols } 0
\DeclareMathSymbol \c_pazo_mapstochar \mathrel { pazosymbols } { "37 }
\DeclareMathSymbol \c_pazo_rightarrow \mathrel { pazosymbols } { "21 }
\cs_new_protected_nopar:Nn \pazo_relbar: {
  \mathrel {
    \mathpalette \mathsm@sh \c_pazo_minus
  }
}
\cs_set_protected_nopar:Nn \MT_mapsto_fill: {
  \arrowfill@ { \c_pazo_mapstochar \pazo_relbar: } \pazo_relbar: \c_pazo_rightarrow
}
\ExplSyntaxOff
\makeatother

\newenvironment{popupbox}[1]%
{%
    \metroset{block=fill}%
    \begin{textblock*}{0.9\textwidth}(0.1\textwidth,0.15\textheight)%
        \begin{alertblock}{#1}%
}%
{       \end{alertblock}%
    \end{textblock*}%
}%


% \metroset{numbering=none}

\usetikzlibrary{tikzmark}

\newcounter{tmp}
\newcommand<>\Highlight[1]{%
\stepcounter{tmp}%
\only#2{\begin{tikzpicture}[remember picture,overlay]
\fill[green!60!black,opacity=0.5] 
  ([xshift=-.2em,yshift=2ex]pic cs:start-\thetmp)
    rectangle  
  ([xshift=.2em,yshift=-1ex]pic cs:end-\thetmp);
\end{tikzpicture}}%
\tikzmark{start-\thetmp}#1\tikzmark{end-\thetmp}%
}

\usepackage{fontawesome5}
\setbeamertemplate{navigation symbols}{}
\AddToHook{shipout/before}{\xdef\currentPageNumber{\inteval{\ReadonlyShipoutCounter+1}}}
\AddToHook{shipout/foreground}{%
  \put(0,0){%
    \raisebox{-\dimexpr\height+0.5ex\relax}[0pt][0pt]{\makebox[\paperwidth][r]{%
      \normalsize\color{structure!40!}%
      \ifnum\currentPageNumber>1%
        \href{\jobname-\zeroPad{\PreviousTotalPages}{\the\numexpr\currentPageNumber-1\relax}.svg}{\faArrowLeft}%
      \else%
        \textcolor{lightgray}{\faArrowLeft}%
      \fi\hspace{0.5ex}%
      \ifnum\currentPageNumber<\PreviousTotalPages%
        \href{\jobname-\zeroPad{\PreviousTotalPages}{\the\numexpr\currentPageNumber+1\relax}.svg}{\faArrowRight}%
      \else%
        \textcolor{lightgray}{\faArrowRight}%
      \fi%
      \hspace{0.5ex}%
    }}%
  }%
  \special{dvisvgm:raw
    <style>svg{cursor: none}</style>%
    <defs>%
    <script type="text/javascript">%
    <![CDATA[%
      var cursorTimer;%
      var downOnLink=false;%
      var downOnRoot=false;%
      function islink(tg){return (tg.tagName=='a') ? true : tg.parentNode ? islink(tg.parentNode) : false;};%
      function ismmedia(tg){return (tg.tagName=='video'||tg.tagName=='audio') ? true : tg.parentNode ? ismmedia(tg.parentNode) : false;};%
      document.addEventListener('mousemove', function(e){%
        if(islink(e.target)||ismmedia(e.target)||e.target.getAttribute('class')=='annot'){%
          e.target.style.cursor='pointer';}else{e.target.style.cursor='default';}%
        try{clearTimeout(cursorTimer);}catch(err){};%
        cursorTimer=setTimeout(function(){e.target.style.cursor='none';},3000);%
      });%
      window.addEventListener('contextmenu', function(e){% capture right click
        if(!islink(e.target)&&!ismmedia(e.target)) e.preventDefault();%
      });%
      document.addEventListener('mousedown', function(e){%
        if(islink(e.target)||ismmedia(e.target)) downOnLink=true;%
        else downOnRoot=true;%
      });%
      document.addEventListener('mouseup', function(e){%
        if(downOnLink||!downOnRoot){downOnLink=false;return;}%
        downOnRoot=false;%
        \ifnum\currentPageNumber<\PreviousTotalPages
          if(!e.shiftKey&&e.button==0) document.location.replace('\jobname-\zeroPad{\PreviousTotalPages}{\the\numexpr\currentPageNumber+1\relax}.svg');%
        \fi%
        \ifnum\currentPageNumber>1
          if(e.shiftKey||e.button>1) document.location.replace('\jobname-\zeroPad{\PreviousTotalPages}{\the\numexpr\currentPageNumber-1\relax}.svg');%
        \fi%
      });%
      document.addEventListener('wheel', function(e){%
        \ifnum\currentPageNumber<\PreviousTotalPages
          if(e.deltaY>0){%
            document.location.replace('\jobname-\zeroPad{\PreviousTotalPages}{\the\numexpr\currentPageNumber+1\relax}.svg');%
          }%
        \fi%
        \ifnum\currentPageNumber>1
          if(e.deltaY<0){%
            document.location.replace('\jobname-\zeroPad{\PreviousTotalPages}{\the\numexpr\currentPageNumber-1\relax}.svg');%
          }%
        \fi%
      });%
      document.addEventListener('keydown',function(e){%
          \ifnum\currentPageNumber<\PreviousTotalPages
            if(e.key=='PageDown'||e.key=='ArrowDown'||e.key=='ArrowRight')%
                document.location.replace('\jobname-\zeroPad{\PreviousTotalPages}{\the\numexpr\currentPageNumber+1\relax}.svg');%
          \fi%
          \ifnum\currentPageNumber>1
            if(e.key=='PageUp'||e.key=='ArrowUp'||e.key=='ArrowLeft')%
              document.location.replace('\jobname-\zeroPad{\PreviousTotalPages}{\the\numexpr\currentPageNumber-1\relax}.svg');%
          \fi%
          if(e.key=='Home')
            document.location.replace('\jobname-\zeroPad{\PreviousTotalPages}{1}.svg');%
          if(e.key=='End')
            document.location.replace('\jobname-\PreviousTotalPages.svg');%
        });%
    ]]>%
    </script>%
    </defs>%
  }%
}%
% helper macro \zeroPad : zero-pads integer according to template,
% e. g. 123 --> 00123 if template is `99999`
% #1: arbitrary integer number as template specifying the
%     width, e. g. `987654' for a width of 6 digits
% #2: the number to be formatted
\def\zeroPad#1#2{\zeroPadI{\zeroTemplate{0}{#1}}{#2}}
%low level macros used by \zeroPad
\def\zeroPadI#1#2{% #1: string of zeros specifying width, #2 number
  \ifnum1#2<1#1
    \zeroPadI{#1}{0#2}%
  \else%
    #2%
  \fi%
}%
\def\zeroTemplate#1#2{% create template (string of zeros) from given num
  \ifnum10#1>1#2
    #1%
  \else%
    \zeroTemplate{0#1}{#2}%
  \fi%
}%


%--------------
% DOCUMENT INFO
%--------------

\title[]{Computational Methods in Epidemic Simulation and Inference}
\subtitle{Not-So-Great Models in Complex Situations}
\author{David Wu}
\date{17 August 2023}
\institute{NUMBAT Seminar}
% \titlegraphic{\vspace{4.5cm}\flushright\includegraphics[width=2cm,height=2cm]{exam_resource/uoa-v-colour2.png}} 
\addbibresource{biblio.bib}
% \addbibresource{biblio_samoa.bib}
% \addbibresource{covid_reports.bib}
% \addbibresource{covid_others.bib}



\begin{document}

\frame[plain,noframenumbering]{\titlepage}


\begin{frame}
  \frametitle{About Me}

  \centering
  \begin{tabular}{m{0.3\textwidth} m{0.3\textwidth} m{0.3\textwidth}}
    \begin{center}\includegraphics[width=0.2\textwidth]{img/uoa-v-colour2.png}\end{center}%
    &\begin{center}\includegraphics[width=0.2\textwidth]{img/python.jpeg}\end{center}%
    &\begin{center}\includegraphics[width=0.2\textwidth]{img/virus.jpg}\end{center}\\%
    \begin{center}{\scriptsize PhD from Dept. Engineering Science at the University of Auckland in New Zealand}\end{center}
    &\begin{center}{\scriptsize Primarily use Python tooling}\end{center}
    &\begin{center}{\scriptsize Research is in the field of infectious diseases}\end{center}
  \end{tabular}
\end{frame}

\begin{frame}
  \frametitle{Overview}
  \begin{enumerate}
    \item \hyperlink{fitting-ode}{Fitting ODE models with generalised profiling}
    \item \hyperlink{network-sim}{Simulating stochastic models on networks (maybe skip)}
    \item \hyperlink{surrogates}{Surrogate modelling techniques for inference}
  \end{enumerate}
\end{frame}

\section{Fitting ODE Models}\label{fitting-ode}

\begin{frame}
  \frametitle{General Problem Setup}

  We have data $y$ at time points $t_i$.

  We assume that there is some underlying state $x$ that we observe to get the data.

  We have some model $\mathcal{M}$ of this underlying state that depends on parameters $\theta$.

  \pause

  We want to:

  \begin{enumerate}
    \item Predict the underlying state $x$ or potential data $y$ at future time points
    \item Determine likely model parameters $\theta$ that could produce the data $y$.
  \end{enumerate}



\end{frame}

\begin{frame}

  \frametitle{ODE Models}
  \only<1>{
    \centering
    \includegraphics[width=0.85\textwidth]{img/sir_diagram.png}
  }
  \only<2>{
    \begin{equation*}
      \begin{aligned}
        \frac{dS}{dt} &= -\beta \frac{SI}{N},\\
        \frac{dI}{dt} &= \beta \frac{SI}{N} - \alpha I,\\
        \frac{dR}{dt} &= \alpha I.
      \end{aligned}
    \end{equation*}
  }
  \only<3>{
    \begin{equation*}
      \frac{d}{dt}\begin{bmatrix}
        S\\I\\R
      \end{bmatrix} = \begin{bmatrix}
        -\beta\frac{SI}{N}\\
        \beta\frac{SI}{N} - \alpha I\\
        \alpha I
      \end{bmatrix}.
    \end{equation*}
  }
  \only<4>{
    \begin{equation*}
      \begin{gathered}
        \frac{d}{dt}(x) = \overbrace{f(x; \theta)}^{\text{RHS}}\\
        x = \begin{bmatrix}
          S & I & R
        \end{bmatrix}^T\\
        \theta = \begin{bmatrix}
          \beta & \alpha
        \end{bmatrix}^T\\
      \end{gathered}
    \end{equation*}
  }

\end{frame}

\begin{frame}[t]
  \frametitle{Solving ODE Models}

  Euler's Method:
  \begin{itemize}
    \item $\frac{dx}{dt} \approx \frac{x(t+\Delta t) - x(t)}{\Delta t}$
    
    \item $x(t + \Delta t) = x(t) + (\Delta t) f(x(t); \theta)$
  \end{itemize}
  \only<2>{
    \centering
    \includegraphics[height=0.7\textheight]{img/euler_method.png}
  }
  \only<3>{
  There are methods with higher-order accuracy (wrt $\Delta t$) that require evaluation of $f$ at multiple $t$.

  Adaptive methods exist that take smaller steps in regions of time where the behaviour is "faster" than expected, in order to improve accuracy.

  These methods can be very slow in particular regions of parameter space. (*Problem 1)
  }
\end{frame}


\begin{frame}[t]
  \frametitle{Fitting ODEs}

  We have some data $y$, and we want to fit an ODE model to it: find $\theta$ that would likely reconstruct $y$.

  We typically work with a simple additive error model:

  \begin{equation*}
    \begin{gathered}
      y(t) = \overbrace{g(x(t))}^{\mathclap{\text{observation model}}} + \varepsilon,\\
      \frac{d}{dt} x(t) = f(x(t); \theta),\\
      \varepsilon \sim \mathcal{N}(0, \Gamma).
    \end{gathered}
  \end{equation*}
\end{frame}

\begin{frame}[t]
  \frametitle{Fitting ODEs: Classical Methods}
  Non-linear Least Squares (frequentist)

  \begin{equation*}
    \min_{\theta, x(0)} \qquad \left\lVert y - g\left(\int_0^{t} f(x(t);\theta)\right)  \right\rVert^2_2
  \end{equation*}
  \pause

  Bayesian Analogue

  \begin{equation*}
    p(\theta, x_0 | y) = \underbrace{\exp\left(-\frac{1}{2}\left\lVert y - g\left(\int_0^{t} f(x(t);\theta)\right) \right\rVert^2_{\Gamma^{-1}}\right)}_{\text{likelihood: } \cdot p(y|\theta, x_0)} \;\cdot \;p(\theta, x_0)
  \end{equation*}

\end{frame}

\begin{frame}
  \frametitle{Problems}

  \begin{enumerate}
    \setlength{\itemsep}{2ex}
    \item Numerical integration can be slow.
    \item Solving the ODE implicitly assumes that the model is correct.
  \end{enumerate}
\end{frame}

\begin{frame}[t]
  \frametitle{A Partial Solution: Collocation on Splines}

  We can bypass the expense of the numerical integration by \textit{collocating} the vector field of a proposed solution against the RHS of the ODE.

  \begin{center}
    \animategraphics[autoplay,loop,palindrome,height=0.7\textheight]{2}{img/quiv/quiver_}{0}{6}
  \end{center}

\end{frame}

\begin{frame}[t]
  \frametitle{A Partial Solution: Collocation on Splines}

  Evaluating this on the entire domain is computationally intractable, so we evaluate at specific ``special'' points.
  The obvious choice is to evaluate this at our data points. But the data is noisy, so taking numerical derivatives is hazardous.

  \pause

  We can solve this using a smoothing spline:

  \begin{enumerate}
    \item Smooth the data $y(t)$ to $\hat{y}(t)$ using splines
    \only<2>{\begin{equation*}
      \begin{gathered}
        \hat{y}(t) = \underbrace{\Phi}_{\mathclap{\text{Spline basis}}} \cdot \overbrace{\hat{c}}^{\mathclap{\text{Spline coefs}}}\\
        \hat{c} = \argmin_c \left\{\left\lVert y - \Phi c \right\rVert^2 + \lambda \left\lVert \frac{d^2}{dt^2}\Phi c \right\rVert^2\right\}
      \end{gathered}
    \end{equation*}}
    \pause
    \item Invert $g$ to get $\hat{x}(t) = g^{-1}(\hat{y}(t))$ 
    \item Fit using the likelihood induced by the loss\begin{equation*}
      \left\lVert \frac{d}{dt} \hat{x} - f(\hat{x}; \theta) \right\rVert^2
    \end{equation*}
  \end{enumerate}

\end{frame}

\begin{frame}
  \centering
  \only<1>{
    \includegraphics[height=0.8\textheight]{img/smoother/smoother_01.png}
  }
  \only<2>{
    \includegraphics[height=0.8\textheight]{img/smoother/smoother_02.png}
  }
  \only<3>{
    \includegraphics[height=0.8\textheight]{img/smoother/smoother_03.png}
  }

 \only<4>{
  \animategraphics[autoplay,loop,palindrome,height=0.8\textheight]{4}{img/smoother/smooth_comp_}{0}{10}
 }
\end{frame}

\begin{frame}
  \frametitle{Problems}
  
  \begin{enumerate}
    \setlength{\itemsep}{2ex}
    \item<1-> Smoothing splines make implicit assumptions about the dynamics of the state
    \only<1>{
      \begin{equation*}
        \hat{c} = \argmin_c \left\{\left\lVert y - \Phi c \right\rVert^2 + \lambda \left\lVert \frac{d^2}{dt^2}\Phi c \right\rVert^2\right\}
      \end{equation*}
    }
    \item<2-> The method falls apart if the observation model $g$ is not invertible
    \begin{itemize}
      \item partial observation of the states of the ODE
      \item missing data at certain time points
      \item identifiability problems
    \end{itemize}
  \end{enumerate}
\end{frame}

\begin{frame}
  \frametitle{Using Auxiliary Information: Generalised Profiling}

  \textcite{ramsay2007} introduce \textit{generalised profiling} to bridge these gaps.

  Core idea: Use the ODE to additionally inform the spline smoothing process.

  \pause

  \begin{equation*}\begin{gathered}
    \hat{\theta} = \argmin_\theta \left\{ \lVert y - \Phi \hat{c}(\theta) \rVert^2 \right\}\\
    \hat{c}(\theta) = \argmin_c \left\{ \lVert y - \Phi c \rVert^2 + \lambda \left\lVert \frac{d}{dt}\Phi c - f(\Phi c; \theta)  \right\rVert^2 \right\}
  \end{gathered}\end{equation*}

\end{frame}

\begin{frame}
  \frametitle{A Likelihood-based Interpretation}
  

\end{frame}


\begin{frame}
  \frametitle{Implementation}

  If we are using splines, we have closed forms for the derivatives.

  We can also use \emph{automatic differentiation} to get gradients of the log-likelihood, so that we can perform optimisation on it.

  We use CasADi for this, and it has interfaces to IPOPT for optimisation (typical domain is optimal control).

  {\begin{center}
  \includegraphics[width=0.4\textwidth]{img/casadi}\end{center}}
  
\end{frame}

\begin{frame}
  \frametitle{A Case Study}

  Measles epidemic in Samoa (Sept 2019 -- Jan 2020)

  \begin{itemize}
    \item nearly 6000 cases (pop. 200 000)
    \item 83 deaths
    \item Anomalously low vaccination coverage \begin{itemize}
      \item 40\% first dose of vaccine
      \item 28\% second dose of vaccine 
    \end{itemize}
  \end{itemize}
  
\end{frame}

\begin{frame}
  \frametitle{Fitting ODEs: Measles in Samoa}
\centering
  \vspace{0.25em}
  \includegraphics[height=0.75\textheight]{img/samoa_data}

  \vspace{-0.25em}
  {\tiny \linespread{0}Plots of the available data for the 2019 Samoan measles outbreak (current hospitalisations and cumulative discharges available, but not shown). Top-left: cumulative reported cases; top-right: cumulative reported deaths; bottom-left: report incidence (daily rate of new cases since last report); bottom-right: cumulative reported hospitalisations.\\}

\end{frame}

% \begin{frame}
%   \frametitle{Fitting ODEs: Measles in Samoa}
%   \centering
%   \includegraphics[height=0.85\textheight]{exam_resource/seirh.png}
% \end{frame}

% \begin{frame}
%   \frametitle{Fitting ODEs: Measles in Samoa}
%   \centering
%   Maximum Likelihood Estimates

%   \includegraphics[width=0.45\textwidth]{img/37_samoa/big_state.pdf}
%   \includegraphics[width=0.45\textwidth]{img/37_samoa/smol_state.pdf}
% \end{frame}

% \begin{frame}
%   \frametitle{Fitting ODEs: Measles in Samoa}
%   \centering
%   \includegraphics[height=0.85\textheight]{img/37_samoa/bivariate_profile.pdf}
% \end{frame}

\section{Network Models}\label{network-sim}

% \begin{frame}
%   \frametitle{Network Models: Summary}

%   Objective: Incorporate heterogeneity and stochasticity into an epidemic model through a contact network, and do it at scale. Strongly driven by modelling COVID-19 in New Zealand.

%   Challenges:
%   \begin{itemize}
%     \item Large number of individuals and possible reactions -- algorithms run into memory or computational limitations
%   \end{itemize}

%   Contributions:
%   \begin{itemize}
%     \item Novel extensions to the Gillespie algorithm to allow for efficient simulation using a rejection-based method.
%     \item Numerous applications of the model for scenario modelling of COVID-19 in New Zealand.
%   \end{itemize}
% \end{frame}

% \begin{frame}
%   \centering
%   \frametitle{Network Models: COVID-19 in NZ}

%   \includegraphics[width=0.8\textwidth]{tikz/covid_timeline_nz}

%   \includegraphics[height=0.6\textheight]{exam_resource/cases_nz_stats_nz}

% \vspace{-1em}
%   {\tiny
%     Count of active cases, recovered individuals and deceased individuals due to COVID-19 in NZ (Source: Stats NZ)}
% \end{frame}


\begin{frame}
  \frametitle{Network Models: The Network}
  \centering
  Bipartite Networks: Explicitly modelling contexts of infection

  \includegraphics[height=0.7\textheight]{img/contact-multiple}

  \hfill {\scriptsize *work of others in the Network Contagion Team}
\end{frame}

\begin{frame}
  \frametitle{Network Models: The Model (roughly)}
  \centering

  \includegraphics[height=0.8\textheight]{img/transmission_diagram_big_no_distancing}
\end{frame}

\begin{frame}
  \frametitle{Network Models: Methods}
  % \centering
  Computational Bottlenecks and Complications:
  \begin{enumerate}[<+->]
    \item Updating rates of neighbours upon becoming infected (exposed)
    \item Separation of probability of (ever) firing and rate of firing
    \item Non-Markovian inter-event/waiting times
  \end{enumerate}
  
  \pause
  Solution: Rejection-based sampling
\end{frame}

% \begin{frame}
%   \frametitle{Network Models: Methods}
%   Example: Updating rates%

%   \only<3->{Instead}%

%   \begin{center}%
%     \includegraphics<1>[height=0.7\textheight]{exam_resource/rej/inf2sus_1}%
%     \includegraphics<2>[height=0.7\textheight]{exam_resource/rej/inf2sus_2}%
%     \includegraphics<3>[height=0.7\textheight]{exam_resource/rej/rej_1}%
%     \includegraphics<4>[height=0.7\textheight]{exam_resource/rej/rej_2}%

%     \vspace{-1em}
%     \only<2>{All nodes must update rates}%
%     \only<3>{Only the infected updates}%
%     \only<4>{Upon an infection event firing, randomly sample and perform acceptance check}%
%   \end{center}%

% \end{frame}

% \begin{frame}
%   \frametitle{Network Models: Methods}
%   Non-Markovian inter-event times

%   \centering
%   \includegraphics[height=0.8\textheight]{img/40_covid/gamma_bounded.pdf}
% \end{frame}

% \begin{frame}
%   \frametitle{Network Models: Results}
%   November 2020 unknown case analysis

%   \centering
%   \includegraphics[height=0.7\textheight]{exam_resource/nov_2020_outbreak.png}
% \end{frame}

% \begin{frame}
%   \frametitle{Network Models: Results}
%   End of August 2021 (Delta) outbreak: AL4 (left) vs AL3 (right)

%   \centering
%   \includegraphics[width=0.475\textwidth]{img/40_covid/results/gp_condition/9_sept_al4.png}
%   \includegraphics[width=0.475\textwidth]{img/40_covid/results/gp_condition/9_sept_al3pess.png}

%   \hfill{\small *uncertainty quantification done by others in the Network Contagion Team}
% \end{frame}

\section{Surrogate Models for Inference}\label{surrogates}

% \begin{frame}
%   \frametitle{Surrogate Modelling}
  
%   Objective: Enable inference of parameters of a complex (network-based) model using a surrogate ODE model.

%   Challenges:
%   \begin{itemize}
%     \item Likelihoods for a stochastic network model are intractable; likelihood-free methods are far too expensive for response planning of outbreaks.
%   \end{itemize}

%   Contribution:
%   \begin{itemize}
%     \item Proof-of-concept study on a simple SIR model.
%   \end{itemize}
% \end{frame}

% \begin{frame}
%   \frametitle{Surrogate Modelling: Methods}
%   \centering
%   \includegraphics<1>[height=0.7\textheight]{exam_resource/Z_S_C}
%   \includegraphics<2>[height=0.7\textheight]{exam_resource/Z_S_E}

%   \only<2>{where $e$ is a multivariate Gaussian}
% \end{frame}

% \begin{frame}
%   \frametitle{Surrogate Modelling: Methods}

%   $\zeta(\theta, \eta) \approx s(\theta) + e(\theta, \eta)$

%   $e(\theta, \eta) \sim \mathcal{N} (\mu(\theta, \eta), \Sigma(\theta, \eta))$

%   To estimate $e$:
%   \begin{enumerate}
%   \item Generate realisations of $\zeta$: $Z = \{\zeta(\theta_1, \eta_1), \zeta(\theta_2, \eta_2), \dots\}$

%   \item Compute $E = \{\zeta(\theta_1, \eta_1) - s(\theta_1), \dots\}$

%   \item Take statistics of $E$: 
%   \begin{itemize}
%     \item $\mu(\theta, \eta) = \frac{1}{|E|}\sum_{e_i \in E} e_i$
%     \item $\Sigma(\theta, \eta) = \frac{1}{|E|^2}\left< e_i - \mu , e_i - \mu\right>$
%   \end{itemize}
%   \end{enumerate}
% \end{frame}

% \begin{frame}
%   \frametitle{Surrogate Modelling: Proof of Concept}
%   Approximate a stochastic network SIR model with an SIR ODE:
%   \centering

%   \includegraphics[height=0.85\textheight]{img/50_surrogate/runrunvar.pdf}
% \end{frame}

% \begin{frame}
%   \frametitle{Surrogate Modelling: Proof of Concept}

%   The Curve Registration Problem

%   \centering
%   $\text{Stoch}(t; \theta, \eta) \approx \text{ODE}(t {\color{red}- \tau}; \theta) + e(\mu, \Sigma)$

%   \includegraphics<1>[height=0.8\textheight]{img/50_surrogate/runrunvar.pdf}
%   \includegraphics<2>[height=0.8\textheight]{img/50_surrogate/k100_state_figure_510}
% \end{frame}

% \begin{frame}
%   \frametitle{Surrogate Modelling: Proof of Concept}
%   A rough linear relationship on log-log between transmission parameter and $\tau$

%   \centering
%   \includegraphics[height=0.85\textheight]{img/50_surrogate/linear_linear_relationship_k}
% \end{frame}

% \begin{frame}{Surrogate Modelling: Proof of Concept}

%   \centering
%   \includegraphics[height=0.85\textheight]{img/50_surrogate/surrogate_r0rec_profile}
% \end{frame}

% \section{Summary}

% \begin{frame}
%   \frametitle{Summary}
%   \begin{enumerate}
%     \item \hyperlink{fitting-ode}{Fitting ODE models with generalised profiling}
%     \item \hyperlink{network-sim}{Simulating stochastic models on networks}
%     \item \hyperlink{surrogates}{Surrogate modelling techniques for inference}
%   \end{enumerate}
% \end{frame}

% \begin{frame}{Thanks}
%   \begin{columns}
%     \begin{column}{0.45\textwidth}
%       \includegraphics[width=\textwidth]{exam_resource/Screenshot from 2023-04-12 16-14-37.png}
%     \end{column}
%     \begin{column}{0.45\textwidth}
%       \includegraphics[width=\textwidth]{exam_resource/XbYFITdo.jpeg}
%     \end{column}
    
%   \end{columns}

% \end{frame}

\begin{frame}[standout]
  \centering
\end{frame}
% 
\end{document}
